% Lösungen Einheit 3
\einheitenkopf{Einheit 3 von 5 · Krümmung und Wendepunkte}
\erg{24}{a) $f''(x) = 6x + 4$, $f'''(x) = 6$ \quad b) $f''(x) = 12x^2 - 2$, $f'''(x) = 24x$ \quad c) $f''(x) = -12x$, $f'''(x) = -12$ \quad d) $f''(x) = 6x^2 + 6x$, $f'''(x) = 12x + 6$}
\erg{25}{a) $f''(x) = 12x - 12 = 0 \Leftrightarrow x = 1$, $f'''(1) = 12 \ne 0$: $W(1 \mid 1)$ \quad b) $f''(x) = 12x^2 - 12 = 0 \Leftrightarrow x = \pm 1$, $f'''(\pm 1) = \pm 24 \ne 0$: $W_1(-1 \mid -4)$, $W_2(1 \mid -4)$ \quad c) linksgekrümmt auf $]{-\infty};\,-1]$ und $[1;\,\infty[$, rechtsgekrümmt auf $[-1;\,1]$ \quad d) $f''(x) = 6x - 12 = 0 \Leftrightarrow x = 2$, $f'''(2) = 6 \ne 0$: $W(2 \mid 1)$; $f'(2) = -1$: Wendetangente $y = -x + 3$ \quad e) $f''(1) = -6 + 6 = 0$, $f'''(1) = -6 \ne 0$, $f(1) = 3$ \quad f) $f'(x) = 12x^3 + 12x^2$, $f''(x) = 36x^2 + 24x$, $f'''(x) = 72x + 24$: $f'(0) = 0$, $f''(0) = 0$, $f'''(0) = 24 \ne 0$, $f(0) = 1$: Wendepunkt mit waagerechter Tangente}
\erg{26}{a) $f'(x) = (x + 1)\,\mathrm{e}^{x}$, $f''(x) = \mathrm{e}^{x} + (x + 1)\,\mathrm{e}^{x} = (x + 2)\,\mathrm{e}^{x}$; $f''(x) = 0 \Leftrightarrow x = -2$, $f'''(-2) = \mathrm{e}^{-2} \ne 0$: $W(-2 \mid -2\mathrm{e}^{-2}) \approx (-2 \mid -0{,}27)$ \quad b) Beim Annähern an die $x$-Achse von unten ist der Graph rechtsgekrümmt, am Tiefpunkt linksgekrümmt; die Krümmung wechselt dazwischen, also gibt es einen Wendepunkt links von $x = 1$. Mögliche Skizze:}

\begin{ksys}[xmin=-6,xmax=4,ymin=-4,ymax=3,klein]
\funktionab{3/2.71828*(\x-2)*exp(\x)}{g}{-6}{2.3}
\tiefpunkt{1}{-3}{T}
\end{ksys}

\erg{27}{a) Steigung $-1$: Steigungswinkel $135^\circ$, die Tangente schließt mit der $x$-Achse einen Winkel von $45^\circ$ ein \quad b) $f''(x) = 12x^2 - 12x = 12x(x - 1)$, $f'''(0) = -12 \ne 0$, $f'''(1) = 12 \ne 0$: $W_1(0 \mid 0)$, $W_2(1 \mid -1)$; Gerade $y = -x$ (durch $(0 \mid 0)$ und $(1 \mid -1)$) \quad c) z.\,B. $y = -x + 3$ (durch $(0 \mid 3)$ und $(3 \mid 0)$); sie trifft den gezeichneten Graphen nur rechts bei $x \approx 2{,}1$. Jede Parallele $y = -x + c$ mit $2 < c \le 4{,}33$ ist richtig.}
\erg{28}{a) Lea vertauscht die Krümmungen: $f'' > 0$ heißt linksgekrümmt. Richtig: linksgekrümmt auf $]{-\infty};\,2]$, rechtsgekrümmt auf $[2;\,\infty[$ \quad b) $g''(x) = 6x + 6$: rechtsgekrümmt auf $]{-\infty};\,-1]$, linksgekrümmt auf $[-1;\,\infty[$}
\erg{29}{a) Ein Wendepunkt ist ein Krümmungswechsel; die Krümmung zeigt das Vorzeichen von $f''$. Nur $f''(x_0) = 0$ reicht nicht: $f(x) = x^4$ hat $f''(0) = 0$ und ist überall linksgekrümmt. \quad b) Auf der linksgekrümmten Seite liegt der Graph oberhalb der Tangente, auf der rechtsgekrümmten unterhalb; im Wendepunkt wechselt er also die Seite. \quad c) $x^3 + x = mx \Leftrightarrow x\,(x^2 + 1 - m) = 0$: für $m > 1$ drei Schnittpunkte, für $m \le 1$ genau einer ($m = 1$ ist die Wendetangente)}
\erg{30}{a) $f''(x) = 0{,}6x - 1{,}2 = 0 \Leftrightarrow x = 2$, $f'''(x) = 0{,}6 \ne 0$, $f(2) = 0{,}4$: Wechsel im Punkt $(2 \mid 0{,}4)$, also 200 m in $x$-Richtung und 40 m in $y$-Richtung \quad b) $f''(x) > 0$ für $x > 2$: der Fahrer lenkt von $x = 200$ m bis $x = 500$ m nach links}
